\section{Auxiliary Motion Objectives}
\label{sec:supp-aux}

The auxiliary objectives supervise frame-to-frame displacement,
cumulative displacement direction, and final displacement,
all in normalized coordinates.

\subsection{Clean-Flow Reconstruction}

We reconstruct clean displacements from the predicted rectified-flow
vector field for auxiliary supervision. Let $\mathbf{x}_0$ denote the
noise-free displacement sequence and $\sigma\in[0,1]$ the noise level.
Under the linear interpolation
$\mathbf{x}_{\sigma}=(1-\sigma)\mathbf{x}_0+\sigma\boldsymbol{\epsilon}$,
the target vector field is
$\mathbf{u}=\boldsymbol{\epsilon}-\mathbf{x}_0$.
The clean-flow estimate is
\begin{equation}
\widehat{\mathbf{x}}_0
=\mathbf{x}_{\sigma}-\sigma\mathbf{v}_{\theta}.
\label{eq:supp-clean-reconstruction}
\end{equation}
The network retains its vector-field parameterization;
auxiliary losses act on the reconstructed sequence.
Here, $x_0$ denotes the noise-free sequence, not the first video frame.

Substituting $\mathbf{x}_{\sigma}=\mathbf{x}_0+\sigma\mathbf{u}$
into Eq.~\eqref{eq:supp-clean-reconstruction} gives
$\widehat{\mathbf{x}}_0-\mathbf{x}_0
=-\sigma(\mathbf{v}_{\theta}-\mathbf{u})$, and hence
\begin{equation}
\|\widehat{\mathbf{x}}_0-\mathbf{x}_0\|_F^2
=\sigma^2\|\mathbf{v}_{\theta}-\mathbf{u}\|_F^2.
\label{eq:supp-clean-error}
\end{equation}
Thus, squared displacement supervision is exactly a noise-dependent
reweighting of the vector-field squared error.

\subsection{Motion Supervision}

Let $\mathbf{f}_{ti}$ and $\widehat{\mathbf{f}}_{ti}$ denote the
ground-truth and reconstructed normalized displacements of point $i$
at step $t$. Their cumulative displacements from the initial position
are $\mathbf{d}_{ti}=\sum_{\tau=1}^{t}\mathbf{f}_{\tau i}$ and
$\widehat{\mathbf{d}}_{ti}
=\sum_{\tau=1}^{t}\widehat{\mathbf{f}}_{\tau i}$.
Each sequence contains $T=48$ displacement steps and $N=64$ points.
We omit the batch index below for readability.

\paragraph{Adjacent-displacement loss.}
The displacement loss penalizes the squared difference between predicted
and ground-truth motion at each step:
\begin{equation}
\ell_{\mathrm{vmag}}(t,i)
=\|\widehat{\mathbf{f}}_{ti}-\mathbf{f}_{ti}\|_2^2.
\label{eq:supp-vmag}
\end{equation}
This loss supervises both magnitude and direction of local 3D
displacements, with its relation to vector-field supervision given
by Eq.~\eqref{eq:supp-clean-error}.

\paragraph{Cumulative-direction loss.}
Small local errors can accumulate into an incorrect overall direction.
We therefore compare the directions of the cumulative displacements:
\begin{equation}
\ell_{\mathrm{dir}}(t,i)
=1-\operatorname{cos}_{\varepsilon}
(\widehat{\mathbf{d}}_{ti},\mathbf{d}_{ti}).
\label{eq:supp-dir}
\end{equation}
Each error is weighted by the ground-truth magnitude
$\|\mathbf{d}_{ti}\|_2$, reducing the contribution of points with little
net displacement. For numerical stability,
$\operatorname{cos}_{\varepsilon}$ computes the dot product divided by
the product of the two vector norms, with each norm clamped below at
$\varepsilon=10^{-6}$.

\paragraph{Relative endpoint loss.}
To constrain the final predicted position, we measure endpoint error
relative to the ground-truth net displacement:
\begin{equation}
\ell_{\mathrm{end}}(i)
=\frac{
\|\widehat{\mathbf{d}}_{Ti}-\mathbf{d}_{Ti}\|_2
}{\max(\|\mathbf{d}_{Ti}\|_2,\delta)},
\qquad \delta=0.1.
\label{eq:supp-end}
\end{equation}
Only points satisfying $\|\mathbf{d}_{Ti}\|_2\geq\delta$ contribute.
This excludes small-net-displacement trajectories, including stationary
motion and motion that returns close to its starting position, for which
relative endpoint error is unstable. The threshold is defined in normalized
coordinates.

\paragraph{Weighting and reduction.}
For sample $b$, all three objectives use the noise weight
$w_b=1-\sigma_b$. The displacement loss averages
$\ell_{\mathrm{vmag}}$ over samples, steps, and points with weights $w_b$.
The direction loss averages $\ell_{\mathrm{dir}}$ over the same indices
with weights $w_b\|\mathbf{d}_{bti}\|_2$.
The endpoint loss averages $\ell_{\mathrm{end}}$ over eligible sample--point
pairs with weights $w_b$. Each weighted sum is divided by its total weight
plus $10^{-6}$; if no endpoint is eligible, its loss is zero.

This weighting gives high-noise samples relatively less influence within
each batch. For the squared displacement term, the coefficient on
vector-field squared error before batch normalization is
$\sigma^2(1-\sigma)$, by Eq.~\eqref{eq:supp-clean-error}.
The direction and endpoint terms impose additional trajectory-level
constraints; they do not share the squared-error equivalence.
Their effects are evaluated through the ablation experiments in the
main paper.

\paragraph{Training configuration.}
For second-stage training, the reduced losses use
$\lambda_{\mathrm{dir}}=0.1$,
$\lambda_{\mathrm{vmag}}=0.1$, and
$\lambda_{\mathrm{end}}=0.05$.